\documentclass[../main.tex]{subfiles}

\begin{document}

\chapter{Kuantor Universal dan Kuantor Eksistensial}

\begin{subcpmk}
  \item Mahasiswa mampu menerjemahkan pernyataan bahasa alami ke dalam ekspresi logika predikat menggunakan kuantor universal dan eksistensial secara lugas.
\end{subcpmk}

\section{Pendahuluan Kuantifikasi Komputasi}
Untuk mengatasi keterbatasan Logika Proposisi dalam merepresentasikan *collection of items* / himpunan instansiasi logik yang sangat besar di ilmu komputer (Database, Array iterasi fungsi), digunakan alat spesifikasi batas kuantitas yang berakar dari operator Kalkulus, yakni \textbf{Kuantor} (Quantifier).

Kuantor men-skala-kan properti "Kebenaran" milik suatu fungsi predikat $P(x)$ agar menjangkau / di\_binding ke setiap elemen domain *Universe Of Discourse* ($U$).

\section{Kuantor Universal ($\forall$)}
Frasa "Semua", "Setiap", "Seluruh", diterjemahkan formal dengan simbol Kuantor Universal ($\forall$, sebuah huruf A yang digambar terbalik dari kata "*All*").

Ekspresinya ditulis matematis $\forall x \ P(x)$, dibaca: \textit{"Untuk semua x (di ranah himpunan $U$), berlaku sifat $P(x)$."}

\begin{konsep}
Kuantor Universal identik hubungannya dengan gerbang **Logika AND ($\land$)** yang diulang masif untuk seluruh anggota $U$.
Bila $U = \{1, 2, 3\}$, maka klaim $\forall x \ P(x)$ secara teknis operasional adalah sintesis evaluasi logikal absolut dari: $P(1) \land P(2) \land P(3)$. Jika ada **Satu Saja** var data yang menyebabkan fungsi P salah (False), maka seluruh pernyataan Kuantor Universal ini sontak runtuh (*False loop breaker*).
\end{konsep}

\section{Kuantor Eksistensial ($\exists$)}
Frasa longgar komputasi pencarian data/ \textit{query filter matching} seperti "Ada", "Terdapat setidaknya satu", atau "Beberapa", diwakili simbol Kuantor Eksistensial ($\exists$, huruf E simetrik dari kata "*Exists*").

Notasi standar komputasinya adalah $\exists x \ P(x)$, secara semantik dinarasikan \textit{"Ada (setidaknya satu) $x$ sedemikian rupa sehingga sifat predikat $P(x)$ itu benar/True."}

\begin{konsep}
Kuantor Eksistensial sejajar dengan gerbang **Logika OR ($\lor$)** makro algoritmik. Bila $U = \{1, 2, 3\}$, klaim $\exists x \ P(x)$ setara ekspresi operasi mesin seleksi komparator Boolean: $P(1) \lor P(2) \lor P(3)$. Cukup kompilator menjumpai **satu saja** temuan data bernilai T, maka perulangan algoritma akan tervalidasi True sukses.
\end{konsep}

\begin{contoh}
Terjemahan Bahasa Alami ke Logika Predikat Komputasi:
- "Setiap \textit{process} komputasi ($P$) yang kelebihan *load* CPU ($L$) akan dibekukan oleh OS server ($B$)".
Terjemahan: $\forall x \ (P(x) \land L(x) \rightarrow B(x))$
- "Ada setidaknya satu \textit{file} ($F$) di \textit{Directory Root} ini yang terinfeksi \textit{malware} ($M$)".
Terjemahan: $\exists x \ (F(x) \land M(x))$
\end{contoh}

\begin{aktivitas}
  \item Berlatihlah di grup menerjemahkan negasi kompleks dari algoritma spesifikasi ini: Apa bedanya algoritma yang dideklarasikan dengan formal $\sim(\forall x P(x))$ dengan yang dikonstruksi secara formal sebagai eksistensial $\exists x (\sim P(x))$? Tunjukkan kesetaraannya secara logikal intuitif!
\end{aktivitas}

\begin{latihan}
  \item Ubah spesifikasi tekstual "Sistem komputasi melarang bahwa \textit{Ada user} yang memiliki kata sandi panjangnya di bawah 5 karakter" menjadi struktur Kuantor lengkap.
\end{latihan}

\begin{checklist}
  \item Paham fungsi pemetaan absolut fungsi AND / OR pada iterasi algoritmis.
  \item Terampil merepresentasikan "Semua" ($\forall$) vs "Sebagian/Terdapat" ($\exists$).
\end{checklist}

\end{document}
